% !TEX root = lectures.tex
\section{Lecture 4}
\subsection{Graph Sparsification, continued}
Recall that last time we wanted to find an $\epsilon$-cut sparsifier for a graph
and we proposed the following algorithm.
\begin{algothm}
    Keep every edge independently with $p = c \frac{\log n}{k \epsilon^2}$
    where $k$ is the size of the min cut in $G$. If you keep an edge,
    make its weight $1/p$. Then this is a $\epsilon$-cut sparsifier.
    \begin{proof*}
        First, let's try a heuristic calculation. Consider
        a cut $(S, \overline{S})$. The weight of edges in the cut
        that are left after the process is a $p \cdot$ a binomial random variable,
        whose expectation is exactly $\frac{|\Cut(S)|}{p}$. Thus, in expectation
        the cut stays roughly the same. The problem is, 
        our concentration bounds don't quite guarantee that this happens for all cuts.
        The probability that we deviate more than a $(1 \pm \epsilon)$ factor,
        by the Chernoff bound, is $2e^{-O(p \Cut_G(S) \epsilon^2)}$. This expression is the largest when
        the cut is the smallest, i.e. the min-cut, but would only give a probability of $1/n^{O(1)}$.
        However, there are exponentially many cuts, so taking a naive union bound will not work.

        We thus use the following fact which will be proved in homework:
        \begin{fact}
            The number of $\alpha$-approximate min cuts is at most $n^{2\alpha}$.
        \end{fact}
        That is, the number of cuts with size at most $\alpha k$ is at most $n^{2\alpha}$.
        Now, we can bucket the cuts by their size and union bound over them individually,
        a so-called \textbf{stratified union bound}.
        Let's have each bucket $i$ contain the cuts of size $[2^{i - 1}k, 2^{i}k)$.
        Now, in this bucket the number edges is outside of $(1 \pm \epsilon)$
        with probability at most $e^{-c 2^{i - 1} \log n}$ and there are only at most $n^{2^{i + 1}}$
        such cuts, so it holds for any of them with probability at most $e^{-c (2^{i - 1} + 2^{i + 1}) \log n} = e^{-O(2^{i - 1} \log n)} = n^{-O(2^{i - 1})}$.
        Now, summing up this probability over the buckets will give a series,
        which is about $\frac{1}{n^{c}} + \frac{1}{n^{2c}} + \frac{1}{n^{4c}} + \dots$, whose sum is bounded by a constant times the first term.
        Thus, with high probability, all cuts are sparsified.
    \end{proof*}
\end{algothm}

\subsection{Linear Programming}
What's the runtime of Gaussian eliminiation?
Some may say that it terminates in $O(n^3)$ steps,
but notice that each step is really $O(1)$ ``field operations.''
If you're over a small finite field, these are fast, but over the reals
you could have to do a really precise operation that takes a long time.
It turns out, Gaussian elimination DOES run in polynomial time
(the number of bits in the input does not blow up too much)
but this is a very nontrival fact to prove, first proved by Bareiss hundreds
of years after the invention of the algorithm.

A natural counterexample to our intuition is the following problem.
Given two polygons, both given as sets of rational numbers that represent their vertices,
find the one with bigger perimeter. The problem is inherently that the square root cannot
be represented accurately without knowing a priori what error margin to round to.
However, one cannot be sure what the error margin is unless you're given the difference in the perimeters to begin with!

Another such miracle is that linear programming over the reals runs in polynomial time.
\begin{definition}
    Let $a_i \in \R^n$, $b_i \in \R$ for $1 \le i \le m$ and $c_j \in \R$ for $1 \le j \le n$.
    Then a linear program is the following minimization problem
    \[ \min_{x_j \in \R} \sum_{j = 1}^{n} c_j x_j : \{ a_i^{\top} x \le b_i \,\,\,\,\, \forall 1 \le i \le m \} \]
    Let $a_i$ be the rows of some matrix $A$ and stack the $b_i$ into a vector $b$.
    Then we can rewrite the constraints as $Ax \le b$.
\end{definition}
It turns out linear programming is a very applicable and extensible framework for solving problems.

\subsection{Basic Applications of Linear Programming}
One simple problem is an allocation setting where you have $n$ jobs and $n$ factories.
Each job is assigned to precisely one factory and assigned job $i$ to factory $j$
has a cost $p_{ij} \ge 0$. You wish to minimize the cost of all jobs in aggregate.
Let $x_{ij}$ be $1$ if job $i$ is assigned to factory $j$ and $0$ if not.
Then the following integer program exactly captures the problem.
\[ \min \sum_{i, j} p_{ij} x_{ij} : \begin{cases} x_{ij} \in \{0, 1\}\\ \sum_{i} x_{ij} = 1 & \forall j \\ \sum_{j} x_{ij} = 1 & \forall i \end{cases} \]
The second and third constraints that every job and every factory have exactly one match.
However, since integer programming is NP-Hard, we instead
\textbf{relax} the integrality constraint to instead just allow $0 \le x_{ij} \le 1$.
This means the set of feasible solutions is larger, so the minimization value of the new LP could be smaller.

For this LP, the problem is now the solution we get may be fractional, where the $x_{ij}$ are
not integers, and so we would not immediately get an allocation. However, it is possible to prove
that the optimal value achieved by this LP is actually EQUAL to the integer program's value.
This is not usually true, we actually got lucky (something called \textbf{total unimodularity} was present here).

Let's look at a different problem where LPs won't immediately get you the optimal solution.
Consider the problem of maximum weighted vertex cover.
\begin{definition}
    The maximum weighted vertex cover (MWVC) problem
    is a problem where given a graph $G = (V, E)$ and a function $w: V \to \R_+$,
    we wish to find a set of vertices $S$ such that all edges are touched by at least one vertex in $S$
    where $\sum_{v \in S} w(v)$ is minimized.
\end{definition}
Now, this problem is NP-Hard. Nevertheless, we can write an IP that represents the situation
and relax it into an LP:
\[ \min \sum_i w(i) x_i : \begin{cases}
    0 \le x_i \le 1 & \forall i \\
    1 \le x_i + x_j & \forall (i, j) \in E
\end{cases} \]
where $x_i$ is meant to be $1$ if we take vertex $i$ into the cover and $0$ if not,
and the second constraint makes sure that we cover all the edges.
Now, given a fractional solution, we will instead \textbf{round} the
answer to get a set of vertices to pick. In this particular instance, we will take vertex $i$
if $x_i \ge 0.5$, otherwise we will not. This will indeed cover every edge; for each edge
$(i,j)$, $x_i + x_j \ge 1$ so one of them is at least $0.5$.
However, we might take more vertices than we need to; indeed we claim that our vertex cover
has weight $2 \cdot OPT$, making it a $2$-approximation.
To see this, notice that for the LP solution $\sum_i w(i) x_i \le OPT$.
However, each $x_i$ at most doubles (ones rounded to $1$ are at most doubled
and ones rounded to $0$ are zeroed), so therefore the returned cover has weight at most $2(\sum_i w(i) x_i) \le 2 \cdot OPT$
as claimed.

Assuming the complexity theory conjecture that Unique Games is NP-Hard,
then this LP is actually the best approximation we can do for the problem!
So far nothing better is known.